% ==========================================================================
%  Chapter 92 -- 16-Storey Seismic XFEM-FE^2 Results (Plan v2 §Fase 6)
% ==========================================================================

\chapter{16-Storey Seismic XFEM-FE\(^2\) Results}
\label{ch:16storey-seismic-xfem-fe2-results}

\paragraph{Normative status.}
This chapter records the staging closure for the L-shaped 16-storey
RC building under multiscale XFEM-FE\(^2\) coupling
(Plan~v2 \S Fase~6). The heavy dynamic run on the Tohoku MYG004
NS+EW+UD scale=1.0 record sits behind
\texttt{main\_lshaped\_multiscale\_16storey} and is
\texttt{scoped\_deferred} per
\path{data/output/validation_reboot/audit_phase6_plan_v2.json}; the
present chapter documents the \emph{staging} evidence that locks the
activation policy and the H1/H2/H3 hypotheses ahead of the heavy
solve.

\section{Configuration (frozen)}

\begin{itemize}
  \item Three column grades by elevation:
        50/35~MPa (storeys 1--6),
        40/28~MPa (storeys 7--12),
        30/21~MPa (storeys 13--16).
  \item Eight embedded rebars per column (Cap.~80) with cover
        positioning preserved from
        \href{run:/memories/repo/lshaped-multiscale-status.md}{%
        \texttt{lshaped-multiscale-status.md}}.
  \item Generalised-\(\alpha\) integrator with \(\rho_{\infty}=0.9\),
        \(\Delta t=0.02\)~s, Rayleigh damping 5\%.
  \item Coupling activated at step~10. Top-5 columns by damage index
        activate XFEM sites at \(\mathrm{NZ}=4\) (Cap.~89 baseline).
\end{itemize}

\section{Staging driver and policy ctest}

\href{run:main_lshaped_multiscale_16storey_staging.cpp}{%
\texttt{main\_lshaped\_multiscale\_16storey\_staging}} is a lightweight
Eigen-only driver that synthesises a 16-storey triangular drift demand
profile peaking at the L-shape transition zones (storeys
\(\sim 5\) and \(\sim 11\)), runs each storey through the 4A planner
+ \texttt{EnrichmentActivationProbe}, and emits the manifest
\path{data/output/validation_reboot/lshaped_multiscale_16storey_staging.json}
with hypotheses H1/H2/H3.

The dedicated regression test
\href{run:tests/test_lshaped_multiscale_16storey_staging_policy_smoke.cpp}{%
\texttt{test\_lshaped\_multiscale\_16storey\_staging\_policy\_smoke}}
locks the same policy in memory and asserts:

\begin{itemize}
  \item \textbf{H1} (localisation at storeys 5--6 / 11--12):
        the top-5 sites by damage index include at least one storey
        within \(\pm 1\) of the expected peaks. PASS.
  \item \textbf{H2} (cyclic accumulation):
        all top-5 sites accumulate damage \(>0.5\) under the
        five-cycle ramped \(50\to250\)~mm protocol. PASS.
  \item \textbf{H3} (staggered convergence):
        every site converges below 6 staggered iterations under the
        canonical Cap.~79 (0.03,\,6) gate. PASS.
\end{itemize}

\section{Scoped-deferred heavy run}

The real dynamic 16-storey solve under MYG004 NS+EW+UD scale=1.0
remains \texttt{scoped\_deferred}. Its runtime is hours on
server-class hardware and requires \texttt{LIBS FULL} with PETSc/SNES
+ \texttt{NonlinearSubModelEvolver} + the seismic harness. The
hypotheses H1/H2/H3 are pre-locked by the staging policy ctest, so
that the heavy run -- when scheduled -- can be evaluated against an
already-frozen acceptance contract.

\paragraph{Honest scientific status.}
\texttt{synthetic\_staging\_no\_real\_dynamic\_solve} for the present
chapter; \texttt{scoped\_deferred\_to\_branch} for the heavy MYG004
run.
